% paper/sections/05_time_integration.tex
% §5：時間積分スキーム — 項別設計と ch13 production stack
%
% ─────────────────────────────────────────────────────────────────────────────
\section{時間積分スキーム — 項別設計と ch13 production stack}
\label{sec:time_int}
% ─────────────────────────────────────────────────────────────────────────────

\noindent\textbf{Part 1 連続定式化との関係：}
本章は，第~\ref{sec:onefluid}章で導出した一流体 NS 方程式
（移流・粘性・圧力・界面張力・浮力の各項；式~\eqref{eq:NS_onefluid}）と
第~\ref{sec:cls_advection}節で確立した CLS 保存形移流方程式を，
第~\ref{sec:CCD}章の CCD/FCCD/UCCD6 空間離散化と整合する形で
1 物理時間ステップに組み立てる手続きを与える．
ch13 production stack（rising bubble + capillary）が採用する
時間積分の項別配分は次表である：

\begin{tcolorbox}[colback=blue!3!white,colframe=blue!50!black,
                  title=\textbf{ch13 production stack：項別時間積分},
                  fonttitle=\small\bfseries,boxrule=0.4pt,arc=2pt]
\small
\renewcommand{\arraystretch}{1.15}
\begin{tabular}{@{}p{0.30\linewidth}p{0.40\linewidth}p{0.22\linewidth}@{}}
\hline
物理項 & 時間積分 & 本章の節 \\
\hline
界面（CLS）移流 & TVD-RK3（SSPRK3） & §\ref{sec:time_level1} \\
NS 対流 & IMEX-BDF2（EXT2 外挿） & §\ref{sec:time_level1} \\
粘性 & implicit-BDF2（GMRES） & §\ref{sec:time_viscous} \\
界面張力 & PPE 内蔵 jump-decomposed CSF & §\ref{sec:time_csf} \\
圧力投影 & IPC + FCCD PPE & §\ref{sec:ipc_derivation} \\
浮力（rising bubble） & balanced-buoyancy（陽的） & §\ref{sec:time_buoyancy} \\
\hline
\end{tabular}
\end{tcolorbox}

本章の組織化原則は，各物理項を独立な剛性源と見なし，
項ごとに時間処理を選択することにある．具体的には：

\begin{itemize}[leftmargin=2em]
  \item \textbf{移流項}：NS 対流に IMEX-BDF2（EXT2 外挿），CLS 界面に TVD-RK3．
        対流 CFL 制約のみを残す（§\ref{sec:time_level1}）．
  \item \textbf{粘性項}：implicit-BDF2 + Jacobi 前処理 GMRES．
        粘性 CFL 制約を消滅させる（§\ref{sec:time_viscous}）．
  \item \textbf{界面張力項}：PPE 右辺の圧力ジャンプ条件
        $[\bnabla p\cdot\hat{\bm n}]_\Gamma=\sigma\kappa\delta_\Gamma$ を陰的に課す
        jump-decomposed CSF．毛管波 BKZ 制約のみが残る（§\ref{sec:time_csf}）．
  \item \textbf{圧力投影}：IPC により $\delta p = p^{n+1} - p^n$ を $\Ord{\Delta t^2}$ で更新し，
        FCCD PPE で同一面投影を実現（§\ref{sec:ipc_derivation}）．
  \item \textbf{浮力項（rising bubble 限定）}：balanced-buoyancy 分解により
        hydrostatic 成分を PPE に吸収し，残差成分のみを陽的に積分
        （§\ref{sec:time_buoyancy}）．
\end{itemize}

本章の節構成は：
精度整合性の原則（§\ref{sec:time_accuracy_principle}），
演算子別剛性マップ（§\ref{sec:time_operator_stiffness}），
移流（§\ref{sec:time_level1}）→粘性（§\ref{sec:time_viscous}）
→界面張力（§\ref{sec:time_csf}）→圧力投影（§\ref{sec:ipc_derivation}）
→浮力（§\ref{sec:time_buoyancy}），
全体時間精度（§\ref{sec:time_accuracy_table}），
時間刻みガイド（§\ref{sec:time_guide}），
他 Level 設定の位置づけ（§\ref{sec:time_level_legacy}）の順である．

% -------------------------------------------------------------------------------
\subsection{精度整合性の原則}
\label{sec:time_accuracy_principle}
% -------------------------------------------------------------------------------

空間離散化が $\Ord{h^p}$ 精度を持つとき，
時間離散化誤差 $\Ord{\Delta t^q}$ が $h$--$\Delta t$ 同時縮小で
空間誤差を支配しないためには $q \ge p/\alpha$（$\alpha$ は $\Delta t\sim h^\alpha$ の比例指数）
が必要である．本稿の CCD/FCCD は $\Ord{h^6}$/$\Ord{h^4}$ であるため，
$\alpha = 1$（対流 CFL 支配）で $q \ge 4$ が理想的だが，
プロジェクション法の圧力--速度結合や二相界面の非線形性により
実際には $q = 2$（IMEX-BDF2 + IPC + implicit-BDF2 粘性）で
検証誤差バジェットが閉じる（第~\ref{sec:verification}章参照）．

項別の時間精度設計指針：

\begin{itemize}[leftmargin=2em]
  \item \textbf{CLS 界面}：TVD-RK3 で $\Ord{\Delta t^3}$．
        $\Delta t_{\text{step}}=\min$ CFL に律速されるため，
        単段 RK3 の精度がそのまま全体に反映される．
  \item \textbf{NS 全体}：IMEX-BDF2 + IPC + implicit-BDF2 で $\Ord{\Delta t^2}$．
        圧力--速度結合と射影閉包が 2 次精度の上限を決める．
  \item \textbf{Splitting 整合性}：本実装は Lie splitting（IMEX-BDF2 を主軸とする
        fractional step）を採用し，splitting 誤差は $\Ord{\Delta t^2}$ である．
        各段の時間精度（NS 対流 2 次・粘性 2 次・圧力投影 2 次）と整合する．
\end{itemize}

% -------------------------------------------------------------------------------
\subsection{演算子別剛性マップ}
\label{sec:time_operator_stiffness}
% -------------------------------------------------------------------------------

一流体式 NS 系
\[
  \rho(\psi)\,\partial_t\bu
    = -\rho(\psi)\,\mathcal{C}(\bu,\bu)
    + \mathcal{V}(\bu) - \bnabla p + \bm{f}_\sigma + \rho(\psi)\bg,
  \qquad \bnabla\cdot\bu = 0
\]
を演算子ごとに分解すると，剛性源が互いに独立であることが分かる：

\begin{center}
\renewcommand{\arraystretch}{1.3}
\begin{tabular}{lll}
\hline
演算子 & 剛性源 & ch13 production stack \\
\hline
$\mathcal{C}(\bu,\bu)$（NS 対流） & 対流 CFL $|\bu|\Delta t/h$ & IMEX-BDF2（EXT2） \\
CLS 移流 & 界面 CFL & TVD-RK3 \\
$\mathcal{V}(\bu)$（粘性） & 粘性 CFL $\nu\Delta t/h^2$ & implicit-BDF2（GMRES） \\
$\bm{f}_\sigma$（界面張力） & 毛管波 $\Delta t_\text{cap}$ & PPE 内蔵 jump CSF \\
$\bnabla p$，$\bnabla\cdot\bu=0$ & 非圧縮拘束 & IPC + FCCD PPE \\
$\rho\bg$（浮力, rising bubble） & 浮力波 $\Delta t \lesssim N_{\mathrm{BV}}^{-1}$ & balanced-buoyancy \\
\hline
\end{tabular}
\end{center}

\noindent\textbf{Splitting 整合性：}
fractional step 構成は Lie splitting（順次適用）に分類されるが，
IMEX-BDF2 の predictor + IPC 投影 + implicit-BDF2 粘性更新を
1 物理ステップ内で組み合わせることで全体 $\Ord{\Delta t^2}$ を維持する．
Strang splitting（対称化）は 2 次精度を保証する一般原理だが，
本実装の射影演算子の冪等性と BDF2 の自己整合性により
Lie splitting でも 2 次精度が達成される．

この項別選択が本章の組織化原則である．
旧来の Level 1/2/3 連続スペクトル枠組み（検証 / 本番 / 剛性極限）は
§\ref{sec:time_level_legacy} で位置づけのみ補足する．

% -------------------------------------------------------------------------------
\subsection{移流項：CLS 用 TVD-RK3 + NS 用 IMEX-BDF2}
\label{sec:time_level1}% backward-compat: 旧 §5.3 (Level 1) label を継承
% -------------------------------------------------------------------------------

CLS 界面移流と NS 対流の時間処理は別系統で扱う：
界面位置の高精度追跡には $\Ord{\Delta t^3}$ の TVD-RK3 を，
NS 対流項には projection 法と相性の良い $\Ord{\Delta t^2}$ の IMEX-BDF2 を採用する．
両者とも陽的処理であるため対流 CFL 制約 $|\bu|\Delta t/h \lesssim 1$ のみが残り，
粘性 CFL は §~\ref{sec:time_viscous} の implicit-BDF2 で消滅させる．

\subsubsection*{CLS 界面移流：TVD-RK3 (SSPRK3)}

CLS 移流方程式 \eqref{eq:cls_advection} および
再初期化方程式 \eqref{eq:cls_reinit_iso} の仮想時間 $\tau$ 積分では，
Shu--Osher 型 TVD-RK3 法（= SSPRK3）\cite{ShuOsher1988} を用いる：
\begin{align}
  q^{(1)} &= q^n + \Delta t \, \mathcal{L}(q^n) \notag\\
  q^{(2)} &= \frac{3}{4}q^n + \frac{1}{4}q^{(1)} + \frac{1}{4}\Delta t\,\mathcal{L}(q^{(1)}) \notag\\
  q^{n+1} &= \frac{1}{3}q^n + \frac{2}{3}q^{(2)} + \frac{2}{3}\Delta t\,\mathcal{L}(q^{(2)})
  \label{eq:tvd_rk3}
\end{align}
ここで $\mathcal{L}(q)$ は空間演算子であり，
ch13 production stack の CLS 移流では \textbf{FCCD フラックス発散}
$\mathcal{L}(\psi)=-\bnabla\cdot(\psi\bu)$
（§~\ref{sec:advection_motivation} の \texttt{FCCDLevelSetAdvection}；
空間演算子 $\mathcal{L}$ の具体的な構成は §~\ref{sec:advection} で詳述）を用いる．
旧 DCCD 経路（§~\ref{sec:dissipative_ccd}）は法的に保持される．
運動量移流 $(u,v)$ の既定空間スキームは UCCD6（§~\ref{sec:uccd6_def}）である．

\noindent\textbf{TVD（Total Variation Diminishing）の定義：}
$TV(q) \equiv \sum_i |q_{i+1} - q_i|$ が時間的に増大しない性質を指す．
TVD-RK3 は，空間離散化を前進 Euler で進めた 1 ステップが TVD となる場合に
その性質を時間方向へ継承する SSP 型 Runge--Kutta 法である．
本稿の FCCD は線形フィルタであり空間 TVD 性を持たないため，
TVD-RK3 との組合せだけで全変動減少は保証されない（注意~\ref{warn:adv_clamp}）．
クランプ処置（式~\eqref{eq:psi_clamp}）と組み合わせることで $\psi \in [0,1]$ を維持する．

\begin{tcolorbox}[warnbox, title={TVD-RK3 の適用範囲}]
\label{warn:tvd_rk3_scope}
TVD-RK3 は以下の 2 箇所にのみ適用する（NS 対流には使わない）：
\begin{itemize}
  \item \textbf{CLS 移流方程式}（式~\eqref{eq:cls_advection}）：
        $\pd{\psi}{t} + \bnabla\cdot(\psi\,\bu) = 0$（保存形）
  \item \textbf{CLS 再初期化方程式}（式~\eqref{eq:cls_reinit_iso}）の仮想時間 $\tau$ 積分（移流部のみ）
\end{itemize}
NS 対流は次小節の IMEX-BDF2（EXT2 外挿）が ch13 既定．
SSPRK3 を NS 全体に適用する全陽的構成は精度基準試験
（§~\ref{sec:verification} の時間収束試験）に限定する．
\end{tcolorbox}

\subsubsection*{NS 対流：IMEX-BDF2（EXT2 外挿）}

水--空気級の密度比・非一様格子・長時間実験では NS 対流項に
\textbf{IMEX-BDF2 法（EXT2 外挿）} を採用する．対流項は
\[
  \mathcal{C}^{\mathrm{EXT2}}
  = 2\mathcal{C}(\bu^n,\bu^n)
    - \mathcal{C}(\bu^{n-1},\bu^{n-1})
\]
で陽的外挿し，粘性ブロックは BDF2 の係数 $\gamma=2/3$ を用いて
\begin{equation}
  \frac{3\bu^* - 4\bu^n + \bu^{n-1}}{2\Delta t}
  = -\mathcal{C}^{\mathrm{EXT2}}
    + \mathcal{V}(\bu^*)
    + \bm{a}^{\,n+1}_{\mathrm{res}}
  \label{eq:predictor_imex_bdf2}
\end{equation}
を解く．ここで $\mathcal{C}(\bu,q) = \bu\cdot\bnabla q$（対流演算子），
$\bm{a}^{\,n+1}_{\mathrm{res}}$ は式~\eqref{eq:buoyancy_gradient_residual_split} の
浮力残差を投影と同じ面演算子で評価した加速度であり，hydrostatic な勾配成分は
圧力側へ残す（balanced-buoyancy 分解；§~\ref{sec:time_buoyancy}）．
Helmholtz 形で書けば
\[
  \bigl(I-\gamma\Delta t\,\mathcal{V}\bigr)\bu^*
  =
  \frac{4}{3}\bu^n-\frac{1}{3}\bu^{n-1}
  +\gamma\Delta t\bigl(-\mathcal{C}^{\mathrm{EXT2}}
  +\bm{a}^{\,n+1}_{\mathrm{res}}\bigr).
\]
続く投影補正も同じ $\gamma\Delta t$ を用いる：
\begin{equation}
  \bu_f^{n+1}
  = \bu_f^* - \gamma\Delta t\,A_f G_f(\delta p),\qquad
  D_f\bu_f^{n+1}=0,
  \label{eq:projection_imex_bdf2_dt}
\end{equation}
すなわち PPE は $D_f A_f G_f$，Corrector は同じ $A_fG_f$ を共有する．
界面張力は PPE 右辺のジャンプ条件として陰的に課す（§~\ref{sec:time_csf}）．

\paragraph{安定性：対流 CFL と IMEX-BDF2 安定領域}
EXT2 は対流項のみを陽的に扱うため対流 CFL 条件
$|\bu|\Delta t/h \le C_{\mathrm{adv}}\sim\Ord{1}$ が必要である．
粘性ブロックは BDF2 で陰的に扱うため，BDF2 が A-stable
（Hairer--Wanner~\cite{HairerWanner1996} Sec.~V.1）であることから
粘性 CFL $\nu\Delta t/h^2$ は時間刻み制約から外れる．
EXT2 は AB2 と同じく $\Ord{\Delta t^2}$ の零安定多段法であり，
NS 対流項の固有値（純虚軸付近）に対して実用上の絶対安定領域を持つ．
最終的な律速制約は対流 CFL と毛管波制約 $\Delta t_{\mathrm{cap}}$
（§~\ref{sec:time_csf}）の最小値である．

\paragraph{旧実装：AB2+IPC（Level 2a）}
\phantomsection\label{sec:ab2_ipc}% backward-compat: 旧 §5.4 sec:ab2_ipc
\phantomsection\label{sec:time_level2}% backward-compat: 旧 §5.4 sec:time_level2 (Phase 6 で §5.10 へ移設予定)

中密度比向けの旧 Level 2a 実装では，対流項に
\textbf{2 段 Adams--Bashforth (AB2) 法}を採用し
van Kan~(1986)~\cite{vanKan1986} の \textbf{Incremental Pressure Correction (IPC) 法}
と組み合わせていた．対角粘性は Crank--Nicolson（§~\ref{sec:time_cn} 参照），
表面張力は半陰的線形化~\cite{AlandVoigt2019}．Predictor の離散式は：
\begin{equation}
\begin{split}
  \frac{\rho^{n+1}}{\Delta t}(\bu^* - \bu^n)
  &= -\rho^{n+1}\!\left[\tfrac{3}{2}\mathcal{C}(\bu^n,\bu^n)
    - \tfrac{1}{2}\mathcal{C}(\bu^{n-1},\bu^{n-1})\right] \\
  &\quad + \frac{1}{2Re}\bigl[
      \mathcal{D}_{\mathrm{d}}(\mu^{n+1},\bu^*)
      + \mathcal{D}_{\mathrm{d}}(\mu^n,\bu^n)\bigr]
    + \frac{1}{Re}\mathcal{D}_{\mathrm{c}}(\mu^n,\bu^n)
  - \bnabla p^n
  + \bm{f}^{n+1}
\end{split}
  \label{eq:predictor_ab2_ipc}
\end{equation}
で，対流項の外挿係数 $(3/2,\,-1/2)$ が AB2，
前ステップ圧力 $-\bnabla p^n$ の陽的組み込みが IPC の本質である
（変密度拡張は Guermond--Minev--Shen~\cite{Guermond2006} を参照）．
ch13 production stack では IMEX-BDF2 + IPC + implicit-BDF2 粘性に統一済みであるが，
本式は §~\ref{sec:full_algorithm} の全体アルゴリズム参照および
過去ベンチマーク（§~\ref{sec:verification}）との整合性のために保持する．

\begin{tcolorbox}[warnbox, title={AB2/EXT2 起動問題（$n=0$）}]
\label{warn:ab2_startup}
AB2 外挿および EXT2 外挿は $n \ge 1$ でのみ適用可能である．
初回ステップ ($n=0$) では前進 Euler 法 ($\alpha_0=1,\,\alpha_{-1}=0$) で代替する．
1 ステップだけの前進 Euler 起動は次時刻の値に $\Ord{\Delta t^2}$ の開始誤差を与えるが，
AB2/EXT2 の零安定性 (Dahlquist の根条件) によりこの開始誤差は後続ステップで
増幅されず，全体精度 $\Ord{\Delta t^2}$ と整合する \cite{DahlquistBjorck1974}．
\end{tcolorbox}

\subsection{IPC 法の核心：圧力増分 \texorpdfstring{$\delta p$}{delta p} の求解}% Phase 2: \subsubsection → \subsection 昇格 (Phase 5 で §5.6 として再構成)
\label{sec:ipc_derivation}

IPC 法では，圧力 Poisson 方程式において絶対圧力 $p^{n+1}$ の代わりに
\textbf{圧力増分 $\delta p \equiv p^{n+1} - p^n$} を未知変数として解く
（§~\ref{sec:algorithm} Step 6--7 参照）：
\begin{align*}
  &\text{PPE：}\quad \bnabla\cdot\!\left(\frac{1}{\rho^{n+1}}\bnabla(\delta p)\right)
    = \frac{1}{\Delta t}\bnabla\cdot\bu^*\\
  &\text{補正：}\quad \bu^{n+1} = \bu^* - \frac{\Delta t}{\rho^{n+1}}\bnabla(\delta p),
  \qquad p^{n+1} = p^n + \delta p
\end{align*}
CCD 演算子 $\mathcal{L}^{\rho}_{\mathrm{CCD}}$ は絶対圧力版と\textbf{完全に同一}であり，
求解対象変数が $p^{n+1}$ から $\delta p$ に変わるだけである．
この定式化により，標準 Chorin 法の人工的圧力境界層に起因する
$\Ord{\Delta t}$ スプリッティング誤差が $\Ord{\Delta t^2}$ に低減される
（Guermond et al., \textit{Comput. Methods Appl. Mech. Engrg.} \textbf{195}, 2006~\cite{Guermond2006}）．

\noindent\textbf{自己完結 proof sketch（主要項のみ）：}
速度 splitting $\bu^{n+1}=\bu^*-\Delta t\,\bnabla(\delta p)/\rho^{n+1}$ と
正確な Navier--Stokes 圧力場 $p_e^{n+1}$ を比較する．
Predictor 段階の誤差は AB2 + CCD で $\bu^*-\bu_e^{n+1}=\Ord{\Delta t^2}$．
PPE は $\bnabla\cdot\bigl(\rho^{-1}\bnabla\delta p\bigr)=\bnabla\cdot\bu^*/\Delta t$ を満たし，
両辺に $\Delta t$ を掛けて発散をとると
$\bnabla\cdot\bu^{n+1}=0$（離散非圧縮）．
$\delta p = p^{n+1}-p^n$ を未知数にすることで境界条件
$\partial_n(\delta p)|_{\partial\Omega}=0$ が
\textbf{時刻 $n+1$ と $n$ の差分ノイマン}として閉じ，
Chorin 法の $\partial_n p^{n+1}|_{\partial\Omega}=0$ 強制が誘起する
$\Ord{\Delta t}$ 境界層が消える．
従って $p_e^{n+1}-p^{n+1}=\Ord{\Delta t^2}$．
速度補正そのものは $\Delta t\bnabla(\delta p)/\rho^{n+1}=\Ord{\Delta t}$ 規模の正則項として残る
（$\delta p=p^{n+1}-p^n=\Ord{\Delta t}$，$\bnabla=\Ord{1/h}$，$\Delta t\sim h$ より$\Delta t\cdot\Ord{\Delta t/h}/\rho=\Ord{\Delta t}$）が，
真値との差すなわち\textbf{補正の誤差}は
$\Delta t\bnabla(p_e^{n+1}-p^{n+1})/\rho^{n+1}=\Ord{\Delta t^3}$ となり，
Predictor 段階の $\Ord{\Delta t^2}$ 誤差と合わせて全体
$\bu^{n+1}-\bu_e^{n+1}=\Ord{\Delta t^2}$．
変密度 $\rho^{n+1}$ が時刻 $n+1$ で評価される点が van Kan の等密度版から
の主な拡張である．非圧縮流れでは圧力場の正則性 $dp/dt = \Ord{1}$ が成立するため，
$\delta p = p^{n+1}-p^n = \Ord{\Delta t}$ は時刻方向 1 ステップで保たれる．
変密度下でこの正則性が破綻しないことは Guermond--Minev--Shen~\cite{Guermond2006}
§6.2 (rotational pressure correction) およびそこで参照される変密度系列で確立されており，
本節の精度評価は変密度下でも保たれる（直接的な variable-density IPC $\Ord{\Delta t^2}$ 誤差解析は Pyo--Shen / Guermond--Quartapelle 系列を参照）．

\subsection{Level 2 の全体的な時間精度}% Phase 2: \subsubsection → \subsection 昇格 (Phase 6 で §5.8 として ch13 中心に書換)
\label{sec:time_accuracy_table}

\begin{center}
\renewcommand{\arraystretch}{1.3}
\begin{tabular}{lcc}
\hline
方程式・項 & 時間積分法 & 時間精度 \\
\hline
CLS 移流 & TVD-RK3 & $\Ord{\Delta t^3}$ \\
NS 対流項（Predictor） & AB2（L2a）/ EXT2（L2b） & $\Ord{\Delta t^2}$ \\
NS 粘性項 & CN（L2a）/ implicit-BDF2（L2b） & $\Ord{\Delta t^2}$ \\
NS 粘性クロス項 & L2a 陽的，L2b 陰的ブロック & $\Ord{\Delta t}$--$\Ord{\Delta t^2}$\,$^\dagger$ \\
Projection スプリッティング & IPC（L2a）/ BDF2 射影補正（L2b） & $\Ord{\Delta t^2}$ \\
表面張力 & 半陰的 CSF / pressure-jump PPE & $\Ord{\Delta t^2}$ \\
\hline
\textbf{全体（閉包済み均質流域）} & & $\boldsymbol{\Ord{\Delta t^2}}$ \\
\textbf{全体（未閉包の界面粘性クロス項）} & & $\boldsymbol{\Ord{\Delta t}}$\,$^\dagger$ \\
\hline
\end{tabular}
\end{center}

{\small $^\dagger$ 変粘性界面（$|\bnabla\mu| \neq 0$）ではクロス微分項 $\mathcal{D}_{\mathrm{c}}$ の陽的処理が $\Ord{\Delta t}$ 誤差を導入し，界面近傍での全体時間精度が $\Ord{\Delta t}$ に律速される（注意~\ref{warn:cn_cross_derivative} 参照）．$\mu_l/\mu_g \gg 1$ の場合は AB2 外挿 $[\frac{3}{2}\mathcal{D}_{\mathrm{c}}^n - \frac{1}{2}\mathcal{D}_{\mathrm{c}}^{n-1}]$ の適用を推奨する．}

{\small $^\ddagger$ 非一様格子で\textbf{Mode 2（ALE 補正なしの再生成）}（§\ref{sec:grid_ale}）を採用すると，格子速度 $\bm{v}_\text{mesh}$ 補正の省略により上記時間精度は $\Ord{\Delta t}$ に劣化する．本稿の Level 2 既定は\textbf{Mode 1（時間固定格子）}を仮定する（ALE 補正付き再生成は将来拡張）．}

本スキームの\textbf{均質流域での全体時間精度は $\Ord{\Delta t^2}$（2次精度）}である．
界面近傍（$\mu_l/\mu_g \gg 1$）では上記 $^\dagger$ の注意が適用され，
交差粘性項を単純陽解法に留める限り $\Ord{\Delta t}$ が上限となる．
現行 ch13 YAML は時間固定の interface-fitted 格子を用いるため，
Mode 2（ALE 補正なし再生成）による $\Ord{\Delta t}$ 劣化は発生しない．
水--空気級では Level 2b の implicit-BDF2 粘性ブロックと
同一面分相投影により，劣化要因を「未閉包の界面演算子」に局在化させる．
空間精度（Dissipative CCD 移流：実効 $\Ord{\varepsilon_d^{\mathrm{adv}} h^2}$ フィルタ誤差，
CCD：$\Ord{h^6}$ 微分精度）は，移流フィルタ誤差 $\Ord{\varepsilon_d^{\mathrm{adv}} h^2}$ が律速する
全体空間精度と整合する．CSF モデル誤差 $\Ord{\varepsilon^2}$ が空間精度の律速であり，
その他のコンポーネント（CCD 微分・PPE 勾配・界面追跡）は
現行の連続圧力場（smoothed Heaviside）のもとで $\Ord{h^6}$ を維持する
（定量値は §~\ref{sec:bf_static_droplet} 参照）．

\noindent\textbf{設計上の注意：「CLS は $\Ord{\Delta t^3}$ なのに全体が $\Ord{\Delta t^2}$ な理由」}

CLS 移流に TVD-RK3（$\Ord{\Delta t^3}$）を採用している一方，NS 対流項は
AB2 または EXT2（$\Ord{\Delta t^2}$）にとどめている．
この非対称性は意図的な設計選択であり，以下の理由による：

界面と流体運動は\textbf{疎結合（operator-split）}として扱われる．
CLS 移流（界面輸送）は密度場 $\rho^{n+1}$ と曲率 $\kappa^{n+1}$ を決定する「前処理」であり，
その時間精度が $\Ord{\Delta t^2}$ 以上であり，かつ界面近傍の陽的交差粘性項や
表面張力卓越時の RC 時刻ずれが支配しない場合，全体の流体運動精度は
NS 方程式側の $\Ord{\Delta t^2}$ で律速する：
\[
  \text{全体時間精度} = \min\!\bigl(\underbrace{\Ord{\Delta t^3}}_\text{CLS},\;
    \underbrace{\Ord{\Delta t^2}}_\text{NS (AB2/IMEX--BDF2)}\bigr) = \Ord{\Delta t^2}
\]
\textbf{CLS の $\Ord{\Delta t^3}$ 高精度化の役割：}
時間精度を $\Ord{\Delta t^3}$ に向上させることには貢献しないが，
界面プロファイルの変形が小さく保たれることで
（i）曲率 $\kappa$ の数値精度が向上し寄生流れを抑制，
（ii）再初期化の必要頻度が減少し質量保存誤差が低減する，
という実用上の優位が得られる．

\noindent\textbf{背景：標準 Chorin 法のスプリッティング誤差（なぜ AB2 だけでは不十分か）}

対流項を AB2（2次精度）に変更するだけでは全体精度は $\Ord{\Delta t^2}$ に\textbf{ならない}．
標準的な Predictor--Corrector 分割（Chorin, 1968~\cite{Chorin1968}）は，
圧力境界条件の不整合（人工的圧力境界層）に起因する $\Ord{\Delta t}$ スプリッティング誤差を独立に内包するため，
対流項の高次化だけで律速が解消されない．
IPC 法による予測ステップへの $p^n$ の組み込みは，このスプリッティング誤差を
2次精度へ低減する標準的な手段の一つであり，本稿ではこの方針を採用する．

% AB2 起動問題 tcolorbox は §5.3 末尾 (sec:ab2_ipc legacy paragraph) に移動 (Phase 2)

\subsection{粘性項：implicit-BDF2 (GMRES)}
\label{sec:time_viscous}

ch13 production stack の粘性項処理は \textbf{implicit-BDF2 法 + Jacobi 前処理 GMRES}
である．Predictor 式 \eqref{eq:predictor_imex_bdf2} の粘性ブロック
$\mathcal{V}(\bu^*)$ を完全陰的に扱うことで純粘性 CFL 制約 $\Delta t < O(h^2)$ を
消滅させる．純陽的処理では水--空気級の $h$ 縮小と $\nu$ 比 $\sim 10$ により
$\Delta t$ が対流 CFL より 1--2 桁厳しくなり実用不可となる．

\paragraph{Helmholtz 形主定式}
BDF2 の係数 $\gamma=2/3$ で陰的化した粘性ブロックは
\begin{equation}
  \bigl(I - \gamma\Delta t\,\mathcal{L}_{\nu}\bigr)\bu^* = \mathrm{RHS}^{n,n-1},
  \qquad
  \mathcal{L}_{\nu} = \rho^{-1}\bnabla\cdot\bigl[\mu(\bnabla + \bnabla^T)\bigr]
  \label{eq:helmholtz_implicit_bdf2}
\end{equation}
の Helmholtz 形となる．ここで $\mathrm{RHS}^{n,n-1}$ は式~\eqref{eq:predictor_imex_bdf2}
右辺の対流外挿 $-\mathcal{C}^{\mathrm{EXT2}}$ と浮力残差 $\bm{a}^{n+1}_{\mathrm{res}}$
および BDF2 履歴項 $(4\bu^n - \bu^{n-1})/3$ をまとめたものである．
左辺演算子 $A = I - \gamma\Delta t\,\mathcal{L}_{\nu}$ は対称正定値であり，
ch13 既定では Krylov 部分空間法 (GMRES) で解く．

\paragraph{GMRES 収束性と Jacobi 前処理}
$A$ の対角項は $\rho_{ij}/\Delta t + (\mu_{i\pm 1/2,j})/(Re\,h^2)$ に支配されるため，
\textbf{Jacobi 前処理} (対角逆数によるスカラー除算) は安価かつ効果的である．
変粘性 $\mu_l/\mu_g \sim 10^2$ の二相流でも，界面厚さ $\sim 3h$ の局在性により
$A$ の条件数は $\Ord{10^3}$ にとどまる．restart 80 GMRES は ch13 既定 $\mathrm{cfl}=0.10$
で 500 反復以内に残差 $10^{-8}$ まで収束する．

\begin{tcolorbox}[resultbox, title={implicit-BDF2 GMRES 収束目安 (ch13 設計値)}]
\label{box:gmres_implicit_bdf2}
\renewcommand{\arraystretch}{1.15}
\begin{tabular}{@{}lccc@{}}
\hline
$\mu_l/\mu_g$ & 残差許容値 & GMRES restart & iterations 目安 \\
\hline
$1$（等粘性） & $10^{-8}$ & $30$ & $\lesssim 50$ \\
$10$ & $10^{-8}$ & $50$ & $\lesssim 150$ \\
$10^2$（空気--水） & $10^{-8}$ & $80$ & $\lesssim 500$ \\
$10^3$（高粘度比） & $10^{-6}$ & $100$ & 適応的調整推奨 \\
\hline
\end{tabular}

\smallskip
\textbf{stagnation 検出と recovery：}
GMRES が restart 周期内で残差停滞 ($r_{k+1}/r_k > 0.99$) を示す場合，
$\Delta t$ を半分に減らして再試行する．ch13 既定では発生頻度 $< 1\%$．
\end{tcolorbox}

\paragraph{安定性：A-stable と粘性 CFL 消滅}
BDF2 は A-stable (Hairer--Wanner~\cite{HairerWanner1996} Sec.~V.1) であり，
粘性ブロックを完全陰的に扱うことで純粘性拡散による $\Delta t < O(h^2)$ 制約は消滅する．
最終的な律速時間刻みは対流 CFL と毛管波制約 $\Delta t_{\mathrm{cap}}$
(§~\ref{sec:time_csf}) の最小値である．

\paragraph{クロス微分項の陽的処理}
implicit-BDF2 でも，変粘性応力テンソルのクロス微分項 $\mathcal{D}_{\mathrm{c}}$ は
$\bu^*$-成分間連立を避けるため陽的に評価する．
旧 Crank--Nicolson (CN) 文脈で議論されてきた制約は implicit-BDF2 でも本質的に
同じ形で残る（変粘性界面で粘性 CFL が局所的に復活；
注意~\ref{warn:cn_cross_derivative}・\ref{warn:cross_cfl}）．

\begin{tcolorbox}[warnbox, title={CN 粘性項のクロス微分：$u$ と $v$ の陰的結合と実装上の分離}]
\label{warn:cn_cross_derivative}

変粘性流体の粘性応力 $\bnabla\cdot[\mu(\bnabla\bu + \bnabla^T\bu)]$ の $x$ 成分は：
\begin{align}
  \bigl[\bnabla\cdot(\mu(\bnabla\bu+\bnabla^T\bu))\bigr]_x
  &= \underbrace{2\frac{\partial}{\partial x}\!\left(\mu\frac{\partial u}{\partial x}\right)
     + \frac{\partial}{\partial y}\!\left(\mu\frac{\partial u}{\partial y}\right)}_{\text{対角項 $\mathcal{D}_{\mathrm{d}}(u)$}}
   + \underbrace{\frac{\partial}{\partial y}\!\left(\mu\frac{\partial v}{\partial x}\right)}_{\text{クロス微分項 $\mathcal{D}_{\mathrm{c}}(v)$}}
  \label{eq:viscous_full_x}
\end{align}
クロス微分項 $\mathcal{D}_{\mathrm{c}}(v)$ を CN で $t^{n+1}$ に置くと $u^*$-$v^*$ の連立系（$2N^2$ 未知数）が生じ，ブロック Thomas 法の優位が失われる．
\textbf{標準対処法：クロス項のみ時刻 $n$ で陽的評価し，対角項だけ CN で半陰的処理}する：
\begin{equation*}
  \frac{u^* - u^n}{\Delta t}
  = \frac{1}{2Re}\bigl[\mathcal{D}_{\mathrm{d}}(u^n) + \mathcal{D}_{\mathrm{d}}(u^*)\bigr]
  + \frac{1}{Re}\,\mathcal{D}_{\mathrm{c}}(v^n)
  + \bm{f}_x^{n+1}
\end{equation*}
これにより $u^*$, $v^*$ を独立にブロック Thomas 法で求解できる．

クロス微分項の1次陽的処理は $\Ord{\Delta t}$ 誤差を導入しうる．
この項が小さい等粘性・低粘性比の問題では他の誤差に埋もれるが，
気液界面で $|\bnabla\mu| \neq 0$ かつ $\mu_l/\mu_g \gg 10$ の場合には，
界面近傍での全体時間精度を $\Ord{\Delta t}$ に律速しうる．
水/空気系では特に顕著であり，AB2 外挿
$[\frac{3}{2}\mathcal{D}_{\mathrm{c}}^n - \frac{1}{2}\mathcal{D}_{\mathrm{c}}^{n-1}]$
の適用を推奨する（式~\eqref{eq:predictor_ab2_ipc} の $\mathcal{C}$ 係数と同様の外挿）．
\end{tcolorbox}

\begin{tcolorbox}[warnbox, title={変粘性界面でのクロス微分陽的処理：粘性 CFL 条件の実質的復活}]
\label{warn:cross_cfl}

気液界面近傍では $\mu$ が数格子幅で $\mu_l/\mu_g \approx 10^2$ 変化する．クロス微分 $\mathcal{D}_{\mathrm{c}}(v)$ を陽的処理すると，界面上でのフォン・ノイマン安定性解析から局所 CFL 条件が復活する：
\begin{equation}
  \Delta t \le C_{\mathrm{cross}} \cdot \frac{h^2}{\Delta\mu / \rho}
  \label{eq:cross_cfl}
\end{equation}
ここで $C_{\mathrm{cross}} = O(1)$，$\Delta\mu = \mu_l - \mu_g$．$\mu_l/\mu_g = 100$ の系では等粘性系の陽的粘性 CFL より約 $1/100$ 倍厳しくなる可能性がある．

\begin{center}
\renewcommand{\arraystretch}{1.3}
\begin{tabular}{lll}
\hline
処理方法 & 安定条件 & 適用場面 \\
\hline
対角粘性項 $\mathcal{D}_{\mathrm{d}}$（CN 法） & 線形拡散部は A-stable & 常に適用 \\
クロス微分項 $\mathcal{D}_{\mathrm{c}}$（陽的） & 式~\eqref{eq:cross_cfl} の制約あり & $\mu_l/\mu_g \lesssim 10$ では無害 \\
\hline
\end{tabular}
\end{center}

\textbf{シミュレーションが爆発した場合：}対流 CFL・毛管波制約（式~\eqref{eq:dt_adv},\eqref{eq:dt_sigma}）を確認後，なお爆発する場合はクロス微分項 $\mathcal{D}_{\mathrm{c}}$ が原因の可能性がある．$\Delta t$ を $1/10$ に減らして安定化するか確認することを推奨する．
\end{tcolorbox}

\begin{tcolorbox}[resultbox, title={CN 対角粘性系の行列構造}]
\label{box:cn_tridiagonal}
対角粘性項 $\mathcal{D}_{\mathrm{d}}(u) = 2\partial_x(\mu\partial_x u) + \partial_y(\mu\partial_y u)$
の CN 半陰的処理で，$x$ 方向項 $(1/2Re)\cdot 2\partial_x(\mu\partial_x u^*)$
をLHSに移項すると（$(1/2Re)\cdot 2 = 1/Re$），行 $j$ の連立系は：
\[
  \left[\frac{\rho^{n+1}_{ij}}{\Delta t}
    + \frac{\mu_{i+1/2,j}+\mu_{i-1/2,j}}{Re\,h^2}\right]u^*_{ij}
  - \frac{\mu_{i+1/2,j}}{Re\,h^2}\,u^*_{i+1,j}
  - \frac{\mu_{i-1/2,j}}{Re\,h^2}\,u^*_{i-1,j}
  = \mathrm{RHS}_{ij}
\]
の\textbf{三重対角形}となる（下・上対角係数は負；$D_{d,x}$ の因子 $2$ と CN の $1/2$ が相殺して $1/Re$）：

\medskip
\textbf{$x$ 方向スウィープ（各行 $j$ 独立）：}
$[\mathbf{L}_x^j]\{u^*_{\cdot,j}\} = \{\mathrm{RHS}^j\}$ を Thomas 法で $\mathcal{O}(N_x)$ 求解．
\begin{itemize}
  \item 主対角：$\rho^{n+1}/\Delta t + (\mu_{i+1/2}+\mu_{i-1/2})/(Re\,h^2)$
  \item 上・下対角：$-\mu_{i+1/2}/(Re\,h^2)$，$-\mu_{i-1/2}/(Re\,h^2)$（ともに負符号）
\end{itemize}

\textbf{$y$ 方向スウィープ（各列 $i$ 独立）：}
$y$ 方向の $D_{d,y}(u) = \partial_y(\mu\partial_y u)$（因子 $1$）に対し：
\begin{itemize}
  \item 上・下対角：$-\mu_{j+1/2}/(2\,Re\,h^2)$（$y$ 方向は $D_{d,y}$ の因子 $1$ と $1/2Re$ でそのまま $1/(2Re)$）
  \item \textbf{$x$ 方向（係数 $-\mu/(Re\,h^2)$）と $y$ 方向（係数 $-\mu/(2\,Re\,h^2)$）は
        絶対値で 2 倍の差}があり，
        ADI スウィープの実装では混同しないこと
        （差分の起源：非圧縮性 $\partial_x u + \partial_y v = 0$ により
        strain-rate tensor 形式の $x$ 方向粘性項が $\partial_x(2\mu\partial_x u)$ となり
        因子 2 を生む；$y$ 方向は対角項 $\partial_y(\mu\partial_y u)$ そのままで
        因子 1．境界条件には依存しない代数恒等式である）．
\end{itemize}
Thomas 法で $\mathcal{O}(N_y)$ 求解．

\textbf{境界条件の組み込み：}
壁面（ノースリップ，$u^*=0$）では対応行を単位行に置き換える．

\textbf{$x$-$y$ 分割と全体精度：}
上記の行・列方向スウィープは，対角粘性項 $\mathcal{D}_{\mathrm{d}}$ を
方向ごとに半陰的に扱う実装単位を示している．
$u^*$-$v^*$ の成分間連立を避けるために陽的評価へ回すのは
交差粘性項 $\mathcal{D}_{\mathrm{c}}$ であり，
$y$ 方向の対角拡散項そのものを常に陽的化するという意味ではない
（§~\ref{sec:time_accuracy_table}，注意~\ref{warn:cn_cross_derivative} 参照）．
\end{tcolorbox}

\paragraph{旧実装：Crank--Nicolson 法 (Level 2a)}
\phantomsection\label{sec:time_cn}% backward-compat: 旧 §5.4.3 sec:time_cn (08c L102 ref)

中密度比向けの旧 Level 2a 実装では粘性項に \textbf{Crank--Nicolson (CN) 法}
を採用していた：
\begin{equation}
  \frac{u^{n+1} - u^n}{\Delta t}
  = \frac{1}{2}\bigl[\mathcal{D}(u^n) + \mathcal{D}(u^{n+1})\bigr]
  \label{eq:crank_nicolson}
\end{equation}
ここで $\mathcal{D}$ は粘性拡散演算子．線形・自己随伴な拡散演算子全てを陰的に
扱う場合，CN は時間方向 $\Ord{\Delta t^2}$ かつ A-stable で，陽的拡散 CFL を
課さない．本実装では交差粘性項を陽的に残すため，変粘性界面では本節の同じ制約
（注意~\ref{warn:cn_cross_derivative}・\ref{warn:cross_cfl}）が現れる．行列構造は
$\mathcal{D}_{\mathrm{d}}$ を方向別 ADI スウィープで三重対角化し
（\ref{box:cn_tridiagonal}）Thomas 法で $\mathcal{O}(N)$ 求解できるが，
ch13 production stack ではグローバルな 2D 系を GMRES で直接解く
（implicit-BDF2 + Jacobi 前処理；本節主体）．

ch13 統一済みであるため CN 法は §~\ref{sec:full_algorithm} の全体アルゴリズム参照
および過去ベンチマーク（§~\ref{sec:verify_cn_viscous}）との整合性のために保持する．

% -------------------------------------------------------------------------------
\subsection{\texorpdfstring{Level 3 時間積分：剛性レジーム向け Radau IIA（設計フックのみ，§\ref{sec:pure_fccd_dns} で空間側を詳述）}{Level 3 時間積分：剛性レジーム向け Radau IIA（設計フックのみ，§10.5 で空間側を詳述）}}
\label{sec:time_level3}
% -------------------------------------------------------------------------------

超高密度比・極低 We・薄膜 DNS など
Level 2 の半陰的手法では時間刻みが不経済となる窓では，
Level 3 の完全陰的時間積分に切り替える．
\textbf{ルーティング閾値（密度比・粘度比・Weber 数の 3 条件同時成立）}は
§\ref{sec:level_selection} に集約済みであり，
本節は時間積分側の構成要素に限定して骨格を提示する：

\begin{itemize}[leftmargin=2em]
  \item \textbf{対流・粘性・ハイパー粘性}：5 次 A/L-stable IRK である
        \textbf{Radau IIA}（Hairer--Wanner \cite{HairerWanner1996}）の 3 段法に UCCD6
        （§\ref{sec:uccd6_def}）を埋め込む．
        CN が対角粘性を無条件安定化するのに対し，Radau IIA は
        UCCD6 ハイパー粘性まで含めて完全陰的に扱う．
  \item \textbf{表面張力}：Li 2022 \cite{Li2022} の完全陰的ジャンプ条件，
        または Denner 2024 \cite{Denner2024} の圧力--速度--界面完全結合．
  \item \textbf{連立解法}：ステージ毎の JFNK（Jacobian-Free Newton--Krylov）反復．
        Level 2 より 1 ステップあたり 5--20 倍高コストだが，
        $\Delta t$ を 1--2 桁緩めうるため剛性領域では総コストが逆転する．
\end{itemize}

\begin{remark}[Level 3 実装の範囲]
\label{rem:level3_scope}
本稿の Level 3 \textbf{時間積分}（Radau IIA + JFNK）は\textbf{設計フックのみ}で提示する．
Level 3 の\textbf{空間離散・連成アーキテクチャ}（純 FCCD DNS, Phase 1--4）は
§\ref{sec:pure_fccd_dns} で理論的に完備に提示しており，
本節は時間積分スキームの観点に限定した概観である．
\texttt{ns\_pipeline.py} の predictor/PPE/corrector 抽象は
既に 3 ステップを因数分解しているため，
PPE ソルバを触らずに predictor を Radau IIA 実装へ置換可能である．
\end{remark}

\noindent
Level 選択トリガー（剛性レジームの自動判定，3 条件同時成立による L2→L3 移行）は
§\ref{sec:level_trigger} 表~\ref{tab:level_trigger} に集約．

% -------------------------------------------------------------------------------
\subsection{毛管波時間刻み制約：波解像の観点}
\label{sec:time_capillary}
% -------------------------------------------------------------------------------

表面張力 $\sigma$ が作用する鋭い界面の最小解像波長 $\lambda = 2h$ に対し，
毛管波の分散関係は $\omega_c^2 = \sigma k^3/(\rho_1+\rho_2)$
（$k = 2\pi/\lambda$）．Brackbill--Kothe--Zemach（1992）\cite{Brackbill1992} の
\textbf{安定性 CFL} は
\[
  \Delta t_\sigma^\text{BKZ} \le \sqrt{\frac{(\rho_1+\rho_2)h^3}{4\pi\sigma}}
\]
であり，陽的 CSF 積分の条件を与える．しかし Denner--van~Wachem
（2015, JCP 285 \cite{DennerVanWachem2015}；2022, JCP 449 \cite{DennerVanWachem2022}）は，
この拘束が純粋な安定性問題ではなく，毛管波自身を\textbf{時間的にも解像する}
必要から来る\textbf{波解像拘束}であることを示した．
両者の $\sqrt{2}$ 比は BKZ が純粋に安定限界
（$\omega_c\Delta t \le \Ord{1}$；陽解法の Nyquist 評価で $\pi/2\approx 1.57$）を与えるのに対し，
Denner--van Wachem が\textbf{毛管波の半周期を少なくとも 1 時間ステップ}で解像する
波解像条件を課す物理的要請の違いに由来する（文献値 $\sqrt{2}$ 比はこの差の order-of-magnitude 表現）：
\begin{equation}
  \boxed{\;
  \Delta t_\text{cap} \le C_\text{wave}\sqrt{\frac{(\rho_1+\rho_2)h^3}{2\pi\sigma}},
  \qquad C_\text{wave}\simeq 0.1\text{--}0.3.
  \;}
  \label{eq:dt_cap_denner}
\end{equation}

\textbf{重要な帰結}：

\begin{enumerate}[leftmargin=2em]
  \item \textbf{表面張力の陰的化では毛管時間刻みは緩まない．}
        半陰的 CSF（Aland--Voigt \cite{AlandVoigt2019}；Bänsch \cite{Baensch2001}）は
        Laplace--Beltrami 成分で無条件安定だが，物理的毛管波が
        解に寄与する限り $\Delta t$ は式~\eqref{eq:dt_cap_denner} で律速される．
        陰的処理はスプリアス剛性固有値を除くが，物理を変えない．
  \item \textbf{UCCD6 や BF 設計は別層}．
        UCCD6 は移流を安定化し，BF ペアは静的力残差を消し，
        波解像は $\Delta t_\text{cap}$ を決める．3 つの拘束は独立合成される．
  \item \textbf{実運用コロラリ}：
        水--空気 ch13（$\rho_1/\rho_2 = 10^3$, $N=128$）では
        $\Delta t_\text{cap}$ が $\Delta t_\text{CFL,adv}$ と $\Delta t_\nu$ より
        厳しい．タイトカップリング陰的スキームが有利なのは
        毛管波が動的に寄与しない領域のみ．
\end{enumerate}

本稿の Level 2 実装では $C_\text{wave}=0.1$--$0.3$ を設計範囲とし，
$0.1$ 未満は過剰計算，$0.3$ 超は波形歪みリスクとする．
測定毛管周期 $T_c = 2\pi/\omega_c$ は，§\ref{sec:benchmarks} の
毛管波ケースで validation 対象とする．

% -------------------------------------------------------------------------------
\subsection{時間刻みガイド：3 Level の律速合成}
\label{sec:time_guide}
% -------------------------------------------------------------------------------

Level 2 の合成 Δt：
\[
  \Delta t_\text{L2}
  \approx \min\!\bigl(\Delta t_\text{CFL,adv},\; \Delta t_\text{cap},\; \Delta t_\text{operator}\bigr)
  \ll \Delta t_\nu
\]
粘性 CFL は CN/implicit-BDF2 半陰的化で消え，
対流 CFL，毛管波拘束，および非一様格子・投影閉包に由来する
演算子固有拘束のみ残る．
$\mu_l/\mu_g \gg 10$ の高粘性比では注意~\ref{warn:cross_cfl} の
クロス微分 CFL が界面近傍で復活しうる．

Level 1 の合成 Δt（検証用）：
\[
  \Delta t_\text{L1}
  = \min\!\bigl(\Delta t_\text{CFL,adv},\; \Delta t_\nu,\; \Delta t_\sigma^\text{BKZ}\bigr)
\]

Level 3 の合成 Δt（剛性レジーム）：
\[
  \Delta t_\text{L3} \approx \Delta t_\text{cap}
\]
のみが残り，対流・粘性・ハイパー粘性の CFL は JFNK 反復コストと引き換えに消滅．

\medskip
\noindent\textbf{実装上のデフォルト}：
第~\ref{sec:benchmarks}章の二相流ベンチマークはすべて Level 2 に属するが，
中密度比・短時間検証は AB2 + IPC + CN，水--空気級の毛管波および浮力駆動気泡は
IMEX--BDF2 + implicit-BDF2 粘性 + 同一面 BF/分相投影を標準実装とする．
格子収束検証（§\ref{sec:time_accuracy_table}）のみ Level 1 を用い，
$\Ord{\Delta t^2}$ NS + $\Ord{\Delta t^3}$ CLS の律速関係を明示する．
Level 3 への切替は剛性トリガー条件に基づく
（§\ref{sec:level_selection}）．

以上で時間積分の骨格が確定した：
Level 2 を主軸とし，検証では Level 1，剛性では Level 3 へ切り替える
連続スペクトルとして時間積分を扱う．
次章では，本章で確立した時間積分骨格の上で，CLS 移流・運動量移流・粘性項の
各物理項に対する CCD/DCCD/UCCD6/FCCD の用途別配分を整備する．
